Indeed, we argue by induction on @@M0@@. When @@M1@@, this is clear since @@M2@@ and @@M3@@ is @@M4@@-regular for every @@M5@@. When @@M6@@, we have an exact sequence @@M7@@ where @@M8@@ is a @@M9@@-sheaf. By 1.7 and the induction hypothesis, we see that if @@M10@@, then @@M11@@ and @@M12@@ is @@M13@@-regular. Thus, by 1.4, @@M14@@ is @@M15@@-regular and @@M16@@ for @@M17@@. Therefore, @@M18@@, which is @@M19@@ by 1.8 (i), since @@M20@@. Finally, @@M21@@ is @@M22@@-regular by 1.3 (i); this gives the first assertion. The second follows from it by 1.8 (ii) and 1.10.

\sgaStableTarget{sga6.e13.para.1-12}\textbf{1.12.} Let @@M23@@ be a noetherian scheme and @@M24@@ an @@M25@@-scheme of finite type. Let @@M26@@ be a \uline{family of classes of coherent sheaves on the fibers of} @@M27@@; that is, for every point @@M28@@ and every extension @@M29@@ of @@M30@@, coherent sheaves @@M31@@ on the algebraic scheme @@M32@@ are given, and we say that an @@M33@@ and an @@M34@@ belong to the same class if there exist @@M35@@-homomorphisms from @@M36@@ into a common extension @@M37@@ of @@M38@@ such that @@M39@@ and @@M40@@ are isomorphic on @@M41@@. Then the family @@M42@@ is said to be \uline{bounded by the coherent sheaf} @@M43@@ \uline{on} @@M44@@, where @@M45@@ is of finite type over @@M46@@, if @@M47@@ is contained in the family of classes of the sheaves @@M48@@ with @@M49@@; we say that @@M50@@ is \uline{bounded} if there exist @@M51@@ as above.

Suppose in addition that @@M52@@ is projective and endowed with a (relatively) ample invertible sheaf @@M53@@. Then @@M54@@ is said to be a @@M55@@-\uline{family}, for a sequence of integers @@M56@@, if every class in @@M57@@ is represented by an @@M58@@ with @@M59@@ algebraically closed, which is a @@M60@@-sheaf.

